\documentclass[conference]{IEEEtran}

\usepackage[utf8]{inputenc}
\usepackage[T1]{fontenc}
\usepackage{graphicx}
\usepackage{booktabs}
\usepackage{url}
\usepackage{amsmath}
\usepackage{hyperref}
\usepackage{tikz}
\usetikzlibrary{arrows.meta, positioning, shapes.geometric, fit}

\hypersetup{colorlinks=true, linkcolor=blue, citecolor=blue, urlcolor=blue}

\title{Agent Notch: An Ambient Desktop Agent Overlay with
       Provider-Agnostic Inference and Phone-Mediated Control}

\author{\IEEEauthorblockN{Project Documentation}
\IEEEauthorblockA{Personal AI Systems Research\\
Project Report}}

\begin{document}
\maketitle

% ============================================================
\begin{abstract}
Desktop AI assistants force a choice between attention cost and
availability: full-screen chat interfaces compete with the user's
work, while terminal-bound agents are invisible to non-technical
users and single-cloud deployments collapse under quota exhaustion
or network loss. We present \textbf{Agent Notch}, a
Dynamic-Island-style ambient overlay that stays docked at the top
edge of the Windows desktop, expands on hover, and is controlled
remotely from a phone over WhatsApp. The system routes inference
through a \emph{provider-agnostic fallback chain} --- an
OpenCode-hosted session tier, an OpenAI-compatible cloud tier, a
local Ollama tier, and a deterministic offline tier --- so that a
reply is always produced even when every network path fails. The
overlay implements strict human-in-the-loop gating: a bounded
first-in-first-out queue (capacity eight), one-at-a-time
execution, and per-action permission prompts (\emph{allow},
\emph{always}, \emph{reject}). Surrounding services include a
persisted news digest, cached weather, GitHub repository
watching, media-key control, screenshot capture, and a scheduled
watchdog with missed-run recovery. We describe the architecture,
the inference layer, the safety model, and the evaluation
methodology built on 27 automated test modules and
failure-injection scenarios. Results show continuous reply
availability across provider and network failures, with the
hosted-session tier delivering the best latency and the local
tier preserving availability during total network loss.
\end{abstract}

\begin{IEEEkeywords}
ambient intelligence, desktop agent overlay, LLM inference
fallback, multi-provider inference, human-in-the-loop, WhatsApp
bot, local large language models
\end{IEEEkeywords}

% ============================================================
\section{Introduction \& Background}
% ============================================================

\subsection{The Rise of Always-On AI Agents}

Large language model (LLM) assistants have moved from
occasional query tools to persistent collaborators that draft
code, summarize documents, monitor repositories, and act on the
user's behalf~\cite{react,tooluse}. As this shift happened, the
\emph{interface} became the bottleneck: a chat window occupies
the screen, breaks the user's focus, and is unreachable when the
user steps away from the machine. At the same time, production
assistants depend on a single hosted provider, so a quota limit
or a regional outage silently turns a ``personal assistant'' into
a dead window. Two problems therefore emerge: (i)~attention cost
--- the assistant competes with the work it is supposed to
support; and (ii)~\emph{service dropout} --- quota exhaustion,
network loss, or a crashed host process removes the assistant
entirely.

\subsection{Problem Statement}

The project addresses three concrete deficiencies in existing
desktop agent designs:

\begin{enumerate}
  \item \textbf{Attention intrusion.} Full-screen assistants
  demand a dedicated region of the display; minimal CLI agents
  are invisible and offer no glanceable status.
  \item \textbf{Single-provider dependence.} Assistants bound to
  one hosted session API die the moment the account runs out of
  credits or the network drops --- there is no fallback path and
  therefore no meaningful availability guarantee.
  \item \textbf{No remote, safe control.} Monitoring or steering
  a desktop agent typically requires physical presence or a
  VPN-bound web console; there is no phone-first channel with
  rigorous permission gating for consequential actions.
\end{enumerate}

\subsection{Core Concepts}

Two conceptual pairs frame the system.

\textbf{Cloud versus local inference.} A hosted session API
provides strong models at the cost of availability; an on-device
runtime (e.g., Ollama~\cite{ollama}) inverts that trade-off.
Provider-agnostic inference treats both as interchangeable tiers
behind one interface, so model quality and availability become
independent axes.

\textbf{Chat-mediated agent control.} Messaging platforms are
already the user's always-on surface. Reusing WhatsApp as the
control plane~\cite{whatsapp} gives remote presence, delivery
receipts, and an auditable message history without building a
second application --- provided that execution is strictly queued
and gated by human permission prompts.

\subsection{Contributions}

The main contributions of this work are:

\begin{enumerate}
  \item A click-through ambient overlay with Dexter-style
  geometry (96$\times$28~px collapsed nub, radius 12, shoulder
  flares) that expands on hover, collapses on cursor-away, and is
  hidden only via CSS --- never by destroying the window --- to
  preserve transparent-window painting on Windows.
  \item Phone-mediated WhatsApp control with a bounded FIFO queue
  (capacity eight), strict one-at-a-time execution, and a
  permission model of \emph{allow / always / reject} prompts.
  \item A \emph{provider-agnostic inference layer} with a
  four-tier fallback chain (hosted session $\rightarrow$
  OpenAI-compatible cloud $\rightarrow$ local Ollama
  $\rightarrow$ deterministic offline responder) that keeps the
  assistant answering under quota exhaustion and total network
  loss.
  \item A reliability layer combining a Windows scheduled-task
  watchdog with missed-run recovery, a self-audit tail reader,
  and an automated test suite of 27 modules used for
  failure-injection evaluation.
\end{enumerate}

% ============================================================
\section{Literature Review \& Research Gap}
% ============================================================

\subsection{Foundational Studies}

\textbf{Attention-aware and peripheral interfaces.} Glanceable
and ambient displays established that information can be
communicated at the periphery of attention without stealing
focus~\cite{ambient}. Apple's Dynamic Island~\cite{hig}
popularized a related pattern in mobile UI: a small, persistent
shape that morphs to surface live activities and expands into
interactive surfaces on demand. The present work applies this
form factor to the desktop edge as an agent status and control
surface.

\textbf{LLM agent architectures.} ReAct-style reasoning loops
couple thought with tool invocation~\cite{react}, and
tool-use/function-calling interfaces standardized how models
request actions~\cite{tooluse}. Orchestration frameworks
popularized multi-step agent graphs~\cite{langchain}, while
autonomous agent projects demonstrated long-running task
execution~\cite{autogpt}. These works focus on \emph{what} the
agent does; comparatively little work addresses where the agent
\emph{lives} on the desktop or how it survives provider failure.

\textbf{Chat platforms as operations channels.} Bots on
WhatsApp and Telegram are widely used for notifications and
ops alerting~\cite{whatsapp,telegram}; permission-gated remote
execution over chat remains comparatively rare, and bounded
queueing with per-action human approval is not standard in the
literature.

\textbf{Local inference runtimes.} Ollama and similar
runtimes~\cite{ollama} brought quantized on-device models to
consumer hardware, making a local availability tier practical.
Work on heterogeneous model routing and fallback
chains~\cite{routing} typically targets cost optimization rather
than availability under failure.

\subsection{Comparison of Existing Approaches}

Designs across the literature fall into four families:
(i)~\emph{single-provider hosted} --- strongest models, zero
availability guarantee under quota or outage;
(ii)~\emph{chained hosted} --- rotation across hosted providers,
which still fails with the network;
(iii)~\emph{local-only} --- resilient but weaker; and
(iv)~\emph{notification-only} --- resilient yet incapable of
steering the agent. None of these families provides an ambient
desktop surface \emph{and} a phone-first control channel
\emph{and} cross-tier failover in one system.

\subsection{Research Gap}

The identified gap is threefold. First, teachers, operators, and
everyday users cannot reliably \emph{monitor and steer} a
personal agent remotely: existing tools assume the user is
seated at the machine, mirroring the monitoring gap reported in
remote-learning analytics where instructors could not observe
learner state at a distance. Second, single-provider assistants
exhibit the same fragility the present work targets as service
dropout: when the hosted session is unavailable, the assistant
returns an ``offline'' error instead of degrading gracefully.
Third, no evaluated system combines (a) an unobtrusive
always-on overlay, (b) phone-first remote control with human
permission gating and bounded queueing, and (c)
provider-agnostic inference that continues to answer offline.
Agent Notch is designed and implemented to close exactly this
gap.
